\documentclass[11pt]{article}
\usepackage[appendix]{calc}
\calcappendix{Appendix}

\begin{document}

\begin{guide}
Chapter 10 leans on a few facts that its sections state and use but that need more care than a lecture allows: a bounded nondecreasing sequence converges, limits obey the usual laws, and a sequence squeezed between two sequences with the same limit has that limit. This appendix proves them. Section A starts from one assumption about the real numbers (completeness) and proves the Monotone Convergence Theorem from it. Section B uses the $\varepsilon$--$N$ definition of a limit to prove the limit facts for sequences. Section C adds four facts needed by the sections on series: limits and inequalities, subsequences, and the limit at infinity of a monotone function (which settles when an improper integral of a positive function converges). Section E proves the limits of $n$-th roots (such as $n^{1/n}\to1$) that the Root Test needs. Section F, written for Chapter 11, proves the facts about power series that the section on their properties states: the radius of convergence exists, and a power series can be differentiated and integrated term by term inside its interval of convergence. Section G, also written for Chapter 11, proves the two facts about endpoints that the section on Taylor series uses: Abel's theorem (the sum of a power series at an endpoint where it converges is the limit of the function from inside) and the behavior of the binomial series at $x=\pm1$. Section H, written for Section 11.4, proves that power series can be multiplied and divided by the rules used to compute the series of $e^x\sin x$ and $\tan x$. Nothing here is needed to \emph{use} the facts, only to see why they are true.
\end{guide}

% =====================================================================
\section{Completeness and the Monotone Convergence Theorem}

\begin{defn}{Upper bound, supremum}{sup}
Let $S$ be a nonempty set of real numbers. A number $M$ is an \emph{upper bound} of $S$ if $s\le M$ for every $s$ in $S$; the set $S$ is \emph{bounded above} if it has an upper bound. A number $M^*$ is the \emph{least upper bound} (or \emph{supremum}) of $S$, written $M^*=\sup S$, if
\begin{enumerate}[label=(\roman*)]
  \item $M^*$ is an upper bound of $S$, and
  \item $M^*\le M$ for every upper bound $M$ of $S$.
\end{enumerate}
Lower bounds, \emph{bounded below} and the \emph{greatest lower bound} $\inf S$ are defined in the same way with the inequalities reversed.
\end{defn}

The supremum need not belong to the set. The set $\{x : 0<x<1\}$ has supremum $1$, which is not in the set, while the set $\{x : 0<x\le 1\}$ has supremum $1$, which is.

\begin{remark}[The Completeness Axiom]
Every nonempty set of real numbers that is bounded above has a least upper bound.
\end{remark}

This is an assumption about the real numbers, not a theorem: it cannot be derived from the rules of arithmetic and order, so we take it as the starting point. It is exactly what the rational numbers lack. The set of rational numbers $q$ with $q^2<2$ is bounded above by $2$, but among the rational numbers it has no \emph{least} upper bound. Every proof in this appendix, and the Monotone Convergence Theorem on which Chapter 10 relies, rests on it.

\begin{defn}{Limit of a sequence}{limit}
A sequence $\{a_n\}$ \emph{converges} to the number $L$, written $\lim_{n\to\infty}a_n=L$, if for every $\varepsilon>0$ there is an integer $N$ such that
\[
\abs{a_n-L}<\varepsilon \quad\text{for every } n>N.
\]
A sequence that does not converge \emph{diverges}. We write $a_n\to\infty$ if for every real number $K$ there is an integer $N$ such that $a_n>K$ for every $n>N$ (and $a_n\to-\infty$ if $a_n<K$ for every $n>N$ in the same way); such a sequence diverges.
\end{defn}

\begin{thm}{Approximation property of the supremum}{approx}
Let $S$ be nonempty and bounded above, and let $M=\sup S$. For every $\varepsilon>0$ there is an element $s$ of $S$ with $M-\varepsilon<s\le M$.
\end{thm}
\begin{proof}
Because $M-\varepsilon<M$ and $M$ is the \emph{least} upper bound, the number $M-\varepsilon$ is not an upper bound of $S$. So some element $s$ of $S$ satisfies $s>M-\varepsilon$. Also $s\le M$, because $M$ is an upper bound of $S$.
\end{proof}

\begin{thm}{Archimedean property}{archimedes}
For every real number $x$ there is a positive integer $n$ with $n>x$. Consequently, for every $\varepsilon>0$ there is a positive integer $N$ with $\dfrac1N<\varepsilon$.
\end{thm}
\begin{proof}
Suppose, to the contrary, that some real number $x$ satisfies $n\le x$ for every positive integer $n$. Then the set $\N=\{1,2,3,\dots\}$ is nonempty and bounded above (by $x$), so by the Completeness Axiom it has a supremum $M$. By Theorem~\ref{thm:approx} with $\varepsilon=1$ there is a positive integer $n$ with $n>M-1$. Then $n+1$ is a positive integer and $n+1>M$, which contradicts the fact that $M$ is an upper bound of $\N$. This proves the first statement.

For the second, apply the first statement to $x=1/\varepsilon$: there is a positive integer $N>1/\varepsilon$, and then $1/N<\varepsilon$.
\end{proof}

\begin{thm}{Monotone Convergence Theorem}{mct}
\begin{enumerate}[label=(\alph*)]
  \item If $\{a_n\}$ is nondecreasing and bounded above, then $\{a_n\}$ converges, and its limit is $\sup\{a_n\}$.
  \item If $\{a_n\}$ is nonincreasing and bounded below, then $\{a_n\}$ converges, and its limit is $\inf\{a_n\}$.
\end{enumerate}
In particular, a bounded monotonic sequence converges.
\end{thm}
\begin{proof}
(a) The set $S=\{a_1,a_2,a_3,\dots\}$ is nonempty and bounded above, so $L=\sup S$ exists by the Completeness Axiom. Let $\varepsilon>0$. By Theorem~\ref{thm:approx} there is an index $N$ with $L-\varepsilon<a_N\le L$. For every $n>N$ we have $a_n\ge a_N$ because the sequence is nondecreasing, and $a_n\le L$ because $L$ is an upper bound. Hence
\[
L-\varepsilon<a_N\le a_n\le L, \qquad\text{so}\qquad \abs{a_n-L}<\varepsilon .
\]
Since $\varepsilon$ was arbitrary, $a_n\to L$.

(b) The sequence $b_n=-a_n$ is nondecreasing, and it is bounded above by $-m$ if $m$ is a lower bound of $\{a_n\}$. By part (a), $b_n\to L'$ with $L'=\sup\{b_n\}$. Because $\abs{a_n-(-L')}=\abs{b_n-L'}$, we get $a_n\to-L'$. Finally $-L'=\inf\{a_n\}$: every $b_n\le L'$ gives $a_n\ge-L'$, so $-L'$ is a lower bound; and if $m'$ is any lower bound of $\{a_n\}$, then $-m'$ is an upper bound of $\{b_n\}$, so $L'\le -m'$, that is $m'\le -L'$.
\end{proof}

\begin{thm}{The limit of a monotone sequence respects bounds}{mctbound}
If $\{a_n\}$ is nondecreasing and $a_n\le M$ for every $n$, then $\lim_{n\to\infty}a_n\le M$ (and $\lim a_n\ge a_1$). If $\{a_n\}$ is nonincreasing and $a_n\ge m$ for every $n$, then $\lim_{n\to\infty}a_n\ge m$.
\end{thm}
\begin{proof}
In the nondecreasing case the limit is $L=\sup\{a_n\}$ by Theorem~\ref{thm:mct}. Since $M$ is an upper bound of $\{a_n\}$ and $L$ is the least upper bound, $L\le M$; and $L\ge a_1$ because $L$ is an upper bound. The nonincreasing case is the same argument applied to $-a_n$.
\end{proof}

\begin{exmp}{A supremum that is not a maximum}{supnotmax}
Let $S=\{1-\tfrac1n : n=1,2,3,\dots\}=\{0,\tfrac12,\tfrac23,\tfrac34,\dots\}$. Show that $\sup S=1$ and that $S$ has no largest element.

\textbf{Solution.} \emph{Step 1: $1$ is an upper bound.} Every element is $1-\frac1n$ with $\frac1n>0$, so every element is less than $1$.

\emph{Step 2: nothing smaller is an upper bound.} Let $M<1$. We must find an element of $S$ bigger than $M$. The number $1-M$ is positive, so by the Archimedean property (Theorem~\ref{thm:archimedes}) there is an integer $N$ with $\frac1N<1-M$. Rearranging, $1-\frac1N>M$, so $M$ is not an upper bound. Hence every upper bound is at least $1$, and $\sup S=1$.

\emph{Step 3: no maximum.} A maximum would be an element equal to $\sup S=1$, but every element is strictly less than $1$.

As a check on the answer, the elements get closer and closer to $1$: $\lim_{n\to\infty}\left(1-\tfrac1n\right)=1$, so $\sup S=1$.
%@ A-ex1 | limit(1 - 1/n, n, oo) == 1 | \sup S=1
\end{exmp}

\begin{exmp}{Monotone Convergence in action: repeated square roots}{mctroots}
Let $a_1=\sqrt2$ and $a_{n+1}=\sqrt{2+a_n}$. Show that $\{a_n\}$ converges and find its limit.

\textbf{Plan.} The theorem needs two facts: the sequence is monotonic and it is bounded. Terms: $a_1\approx1.414$, $a_2\approx1.848$, $a_3\approx1.962$. They increase; they seem to stay below $2$.
%@ A-ex2a | abs(sqrt(2+sqrt(2)).evalf() - 1.8478) < 1e-4 | 1.848

\textbf{Bounded above by $2$} (induction). $a_1=\sqrt2<2$. If $a_n<2$, then $a_{n+1}=\sqrt{2+a_n}<\sqrt{2+2}=2$, because the square root is increasing. So $a_n<2$ for all $n$.

\textbf{Increasing} (induction). $a_2>a_1$ since $\sqrt{2+\sqrt2}>\sqrt2$. If $a_n>a_{n-1}$, then $2+a_n>2+a_{n-1}$, and taking square roots, $a_{n+1}>a_n$. So the sequence is increasing.

\textbf{Conclusion.} By the Monotone Convergence Theorem (Theorem~\ref{thm:mct}) the sequence converges, say to $L$. Taking the limit in the recurrence (the limit of $a_{n+1}$ is also $L$, and the square root is continuous) gives $L=\sqrt{2+L}$, so $L^2-L-2=0$, that is $(L-2)(L+1)=0$.
%@ A-ex2b | expand((Symbol('L')-2)*(Symbol('L')+1)) == Symbol('L')**2 - Symbol('L') - 2 | (L-2)(L+1)=0
Because $a_n>0$ for all $n$, the limit is not $-1$. Therefore $L=2$, and by Theorem~\ref{thm:mctbound} this agrees with $L\le 2$.
\end{exmp}

\begin{exmp}{A nonincreasing sequence bounded below}{mctnewton}
Let $a_1=2$ and $a_{n+1}=\dfrac12\left(a_n+\dfrac2{a_n}\right)$. Show that $\{a_n\}$ converges and find its limit.

\textbf{Terms.} $a_2=\frac12\left(2+1\right)=\frac32$ and $a_3=\frac12\left(\frac32+\frac43\right)=\frac{17}{12}$.
%@ A-ex3a | (Q(2)+2/Q(2))/2 == Q(3,2) | \left(2+1\right)=\frac32
%@ A-ex3b | (Q(3,2)+2/Q(3,2))/2 == Q(17,12) | \frac{17}{12}

\textbf{Bounded below by $\sqrt2$.} For $a_n>0$,
\[
a_{n+1}^2-2=\frac14\left(a_n+\frac2{a_n}\right)^2-2=\frac14\left(a_n-\frac2{a_n}\right)^2\ge0 ,
\]
because $\left(u+v\right)^2-4uv=(u-v)^2$ with $u=a_n$, $v=2/a_n$ and $uv=2$. So $a_{n+1}\ge\sqrt2$ for every $n$, and $a_1=2>\sqrt2$ too. All terms are positive, so the argument applies at every step.
%@ A-ex3c | simplify(((n + 2/n)/2)**2 - 2 - (n - 2/n)**2/4) == 0 | \frac14\left(a_n-\frac2{a_n}\right)^2

\textbf{Nonincreasing.} Since $a_n\ge\sqrt2$ we have $a_n^2\ge2$, so $\frac2{a_n}\le a_n$ and
\[
a_{n+1}-a_n=\frac12\left(\frac2{a_n}-a_n\right)\le0 .
\]
%@ A-ex3d | simplify(((n + 2/n)/2 - n) - (2/n - n)/2) == 0 | \frac12\left(\frac2{a_n}-a_n\right)

\textbf{Conclusion.} By the Monotone Convergence Theorem, part (b), the sequence converges to some $L\ge\sqrt2$. Taking limits in the recurrence gives $L=\frac12\left(L+\frac2L\right)$, so $2L=L+\frac2L$, $L=\frac2L$, $L^2=2$. Since $L\ge\sqrt2>0$, $L=\sqrt2$. (This is Newton's method for $\sqrt2$.)
\end{exmp}

% =====================================================================
\section{Limit Facts for Sequences}

Throughout, ``for large $n$'' means ``for every $n$ greater than some integer $N$''. We use the $\varepsilon$--$N$ definition (Definition~\ref{def:limit}) and the triangle inequality $\abs{u+v}\le\abs u+\abs v$.

\begin{thm}{Uniqueness of limits}{unique}
A sequence has at most one limit.
\end{thm}
\begin{proof}
Suppose $a_n\to L$ and $a_n\to L'$ with $L\ne L'$. Let $\varepsilon=\frac12\abs{L-L'}>0$. There are integers $N_1,N_2$ with $\abs{a_n-L}<\varepsilon$ for $n>N_1$ and $\abs{a_n-L'}<\varepsilon$ for $n>N_2$. For $n>\max\{N_1,N_2\}$,
\[
\abs{L-L'}\le\abs{L-a_n}+\abs{a_n-L'}<2\varepsilon=\abs{L-L'},
\]
a contradiction. So $L=L'$.
\end{proof}

\begin{thm}{Convergent sequences are bounded}{bounded}
If $\{a_n\}$ converges, then there is a number $K$ with $\abs{a_n}\le K$ for every $n$.
\end{thm}
\begin{proof}
Let $a_n\to L$. Taking $\varepsilon=1$, there is $N$ with $\abs{a_n-L}<1$ for $n>N$, so $\abs{a_n}\le\abs L+1$ for $n>N$. Take $K=\max\{\abs{a_1},\dots,\abs{a_N},\abs L+1\}$.
\end{proof}

\begin{thm}{Tails and absolute values}{tail}
\begin{enumerate}[label=(\alph*)]
  \item If $a_n=b_n$ for all $n>M$ (for some $M$), then $\{a_n\}$ converges if and only if $\{b_n\}$ converges, and then the limits are equal. Likewise $a_n\to L$ if and only if $a_{n+p}\to L$ for any fixed integer $p\ge0$.
  \item $a_n\to0$ if and only if $\abs{a_n}\to0$.
\end{enumerate}
\end{thm}
\begin{proof}
(a) Let $b_n\to L$ and $\varepsilon>0$. There is $N$ with $\abs{b_n-L}<\varepsilon$ for $n>N$. Then for $n>\max\{N,M\}$ we have $a_n=b_n$, so $\abs{a_n-L}<\varepsilon$. The converse is the same with the roles exchanged. The shifted sequence $a_{n+p}$ agrees, after the index change $m=n+p$, with the terms $a_m$, $m>N+p$, so the same $N$ works with $N+p$ replaced by $N$.

(b) The statements $\abs{a_n-0}<\varepsilon$ and $\abs{\,\abs{a_n}-0\,}<\varepsilon$ are the same inequality.
\end{proof}

\begin{thm}{Limit laws for sequences}{laws}
Let $a_n\to A$ and $b_n\to B$, and let $c$ be a constant. Then
\begin{enumerate}[label=(\arabic*)]
  \item $a_n\pm b_n\to A\pm B$;
  \item $ca_n\to cA$;
  \item $a_nb_n\to AB$;
  \item $a_n/b_n\to A/B$, provided $B\ne0$.
\end{enumerate}
\end{thm}
\begin{proof}
(1) Let $\varepsilon>0$. Choose $N_1,N_2$ with $\abs{a_n-A}<\varepsilon/2$ for $n>N_1$ and $\abs{b_n-B}<\varepsilon/2$ for $n>N_2$. For $n>\max\{N_1,N_2\}$,
\[
\abs{(a_n\pm b_n)-(A\pm B)}\le\abs{a_n-A}+\abs{b_n-B}<\varepsilon .
\]

(2) If $c=0$ the claim is trivial. Otherwise choose $N$ with $\abs{a_n-A}<\varepsilon/\abs c$ for $n>N$; then $\abs{ca_n-cA}=\abs c\,\abs{a_n-A}<\varepsilon$.

(3) By Theorem~\ref{thm:bounded} there is $K>0$ with $\abs{a_n}\le K$ for all $n$. Write
\[
a_nb_n-AB=a_n(b_n-B)+B(a_n-A),
\]
so $\abs{a_nb_n-AB}\le K\abs{b_n-B}+\abs B\,\abs{a_n-A}$. Choose $N_1$ with $\abs{b_n-B}<\frac{\varepsilon}{2K}$ for $n>N_1$ and $N_2$ with $\abs{a_n-A}<\frac{\varepsilon}{2(\abs B+1)}$ for $n>N_2$. For $n>\max\{N_1,N_2\}$ the right side is less than $\frac\varepsilon2+\frac\varepsilon2=\varepsilon$ (using $\abs B<\abs B+1$).

(4) First we show $1/b_n\to1/B$. Since $B\ne0$, take $\varepsilon_0=\abs B/2$: there is $N_0$ with $\abs{b_n-B}<\abs B/2$ for $n>N_0$, hence $\abs{b_n}>\abs B/2$ (so $b_n\ne0$) for $n>N_0$. For such $n$,
\[
\left|\frac1{b_n}-\frac1B\right|=\frac{\abs{B-b_n}}{\abs{b_n}\,\abs B}<\frac{2}{\abs B^{\,2}}\,\abs{b_n-B}.
\]
Given $\varepsilon>0$, choose $N\ge N_0$ with $\abs{b_n-B}<\frac{\varepsilon\abs B^{\,2}}2$ for $n>N$; then the left side is less than $\varepsilon$. Now $a_n/b_n=a_n\cdot(1/b_n)\to A\cdot(1/B)$ by part (3).
\end{proof}

\begin{thm}{Squeeze Theorem for sequences}{squeeze}
Suppose $a_n\le b_n\le c_n$ for every $n$ greater than some index $M$. If $a_n\to L$ and $c_n\to L$, then $b_n\to L$.
\end{thm}
\begin{proof}
Let $\varepsilon>0$. Choose $N_1$ with $\abs{a_n-L}<\varepsilon$ for $n>N_1$, and $N_2$ with $\abs{c_n-L}<\varepsilon$ for $n>N_2$. For $n>\max\{M,N_1,N_2\}$,
\[
L-\varepsilon<a_n\le b_n\le c_n<L+\varepsilon,
\]
so $\abs{b_n-L}<\varepsilon$.
\end{proof}

\begin{thm}{Limits of $f(a_n)$ and of $f(n)$}{fofa}
\begin{enumerate}[label=(\alph*)]
  \item If $a_n\to L$ and $f$ is defined at every $a_n$ and is continuous at $L$, then $f(a_n)\to f(L)$.
  \item If $f(n)=a_n$ for every positive integer $n$ and $\lim_{x\to\infty}f(x)=L$, where $L$ is a real number, $\infty$ or $-\infty$, then $\lim_{n\to\infty}a_n=L$.
\end{enumerate}
\end{thm}
\begin{proof}
(a) Let $\varepsilon>0$. By continuity of $f$ at $L$ there is $\delta>0$ such that $\abs{f(y)-f(L)}<\varepsilon$ whenever $\abs{y-L}<\delta$. Since $a_n\to L$ there is $N$ with $\abs{a_n-L}<\delta$ for $n>N$. For those $n$ we get $\abs{f(a_n)-f(L)}<\varepsilon$.

(b) Suppose first that $L$ is a real number and let $\varepsilon>0$. Because $f(x)\to L$ as $x\to\infty$, there is a number $X$ with $\abs{f(x)-L}<\varepsilon$ for all $x>X$. Let $N$ be an integer with $N\ge X$. For every $n>N$ we have $n>X$, so $\abs{a_n-L}=\abs{f(n)-L}<\varepsilon$. If $L=\infty$ and $K$ is any number, there is $X$ with $f(x)>K$ for $x>X$, and again $N\ge X$ gives $a_n=f(n)>K$ for $n>N$. The case $L=-\infty$ is the same.
\end{proof}

The converse of part (b) is false: if $a_n=\cos(2\pi n)$ then $a_n\to1$, but $\lim_{x\to\infty}\cos(2\pi x)$ does not exist.

\begin{exmp}{Finding $N$ from $\varepsilon$}{findN}
Prove that $\displaystyle\lim_{n\to\infty}\frac{2n+1}{n+3}=2$, and find an $N$ that works for $\varepsilon=0.001$.

\textbf{Plan.} We must make $\abs{a_n-2}$ small, so first simplify $\abs{a_n-2}$ to something we can solve for $n$.

\textbf{Solution.} For $n\ge1$,
\[
\left|\frac{2n+1}{n+3}-2\right|=\left|\frac{2n+1-2n-6}{n+3}\right|=\frac{5}{n+3}.
\]
%@ B-ex1a | simplify((2*n+1)/(n+3) - 2 + 5/(n+3)) == 0 | \frac{5}{n+3}
Now $\frac5{n+3}<\varepsilon$ exactly when $n+3>\frac5\varepsilon$, that is, when $n>\frac5\varepsilon-3$. So given any $\varepsilon>0$, any integer $N\ge\frac5\varepsilon-3$ works: for every $n>N$ we get $n+3>\frac5\varepsilon$ and so $\abs{a_n-2}<\varepsilon$. That proves the limit is $2$.

For $\varepsilon=0.001$: $\frac5\varepsilon-3=4997$, so $N=4997$ works. Check at the first index covered, $n=4998$: $\frac{5}{5001}<0.001$.
%@ B-ex1b | Q(5,1)/Q(1,1000) - 3 == 4997 | so $N=4997$ works
%@ B-ex1c | Q(5,5001) < Q(1,1000) | \frac{5}{5001}<0.001
\end{exmp}

\begin{exmp}{A Squeeze Theorem argument}{squeezeex}
Find $\displaystyle\lim_{n\to\infty}\frac{n\sin n}{n^2+1}$.

\textbf{Solution.} The numerator oscillates, so the limit laws do not apply directly (they need $\lim\sin n$ to exist, and it does not). Instead bound the terms. Since $-1\le\sin n\le1$ and $n^2+1>0$,
\[
-\frac{n}{n^2+1}\le\frac{n\sin n}{n^2+1}\le\frac{n}{n^2+1}.
\]
Both outer sequences tend to $0$ because
\[
\lim_{n\to\infty}\frac{n}{n^2+1}=\lim_{n\to\infty}\frac{1/n}{1+1/n^2}=\frac{0}{1+0}=0
\]
by the limit laws (Theorem~\ref{thm:laws}).
%@ B-ex2 | limit(n/(n**2+1), n, oo) == 0 | \frac{0}{1+0}=0
By the Squeeze Theorem (Theorem~\ref{thm:squeeze}), the limit is $0$.
\end{exmp}

\begin{exmp}{Proving that a sequence diverges}{divproof}
Prove directly from the definition that $a_n=(-1)^n$ diverges.

\textbf{Plan.} If the sequence had a limit $L$, then eventually \emph{all} terms would be within a small distance of $L$, including two consecutive terms. But consecutive terms are $2$ apart, so they cannot both be close to one number.

\textbf{Proof.} Suppose $a_n\to L$. Take $\varepsilon=1$. There is $N$ with $\abs{a_n-L}<1$ for all $n>N$. Pick any $n>N$; then $n$ and $n+1$ are both greater than $N$, and
\[
2=\abs{a_n-a_{n+1}}=\abs{(a_n-L)+(L-a_{n+1})}\le\abs{a_n-L}+\abs{L-a_{n+1}}<1+1=2 .
\]
%@ B-ex3 | abs((-1)**3 - (-1)**4) == 2 | 2=\abs{a_n-a_{n+1}}
This says $2<2$, which is impossible. Hence no limit exists. (The same argument with $\varepsilon$ smaller than half the gap shows that any sequence whose terms keep jumping between two different values diverges.)
\end{exmp}

\begin{exmp}{A limit of the form $f(a_n)$}{fofaex}
Find $\displaystyle\lim_{n\to\infty}\cos\!\left(\frac{\pi n}{2n+1}\right)$.

\textbf{Solution.} Let $a_n=\dfrac{n}{2n+1}$. Dividing numerator and denominator by $n$, $a_n=\dfrac{1}{2+1/n}\to\dfrac12$ by the limit laws.
%@ B-ex4a | limit(n/(2*n+1), n, oo) == Q(1,2) | \dfrac12
The function $f(y)=\cos(\pi y)$ is continuous at every number, in particular at $\frac12$. By Theorem~\ref{thm:fofa}(a),
\[
\cos(\pi a_n)\to\cos\!\left(\frac\pi2\right)=0 .
\]
%@ B-ex4b | limit(cos(pi*n/(2*n+1)), n, oo) == 0 | \cos\!\left(\frac\pi2\right)=0
\end{exmp}

% =====================================================================
\section{Limits, Inequalities and Monotone Functions}

This section supports the study of series. It uses Definition~\ref{def:limit}, the supremum facts of Section~A and the limit facts of Section~B.

\begin{thm}{Limits preserve non-strict inequalities}{ineq}
Suppose $a_n\to A$, $b_n\to B$, and $a_n\le b_n$ for every $n$ greater than some index $M$. Then $A\le B$. In particular, if $a_n\ge m$ for large $n$ and $a_n\to A$, then $A\ge m$.
\end{thm}
\begin{proof}
Suppose, to the contrary, that $A>B$, and let $\varepsilon=\frac12(A-B)>0$. There is $N$ with $a_n>A-\varepsilon=\frac{A+B}2$ and $b_n<B+\varepsilon=\frac{A+B}2$ for all $n>N$. For $n>\max\{M,N\}$ this gives $b_n<\frac{A+B}2<a_n$, contradicting $a_n\le b_n$. So $A\le B$. For the second statement take the constant sequence $b_n=m$, applied with the roles of the two sequences exchanged: $m\le a_n$ gives $m\le A$.
\end{proof}

Strict inequalities are not preserved: $\frac1{n^2}<\frac1n$ for every $n\ge2$, yet both sequences tend to $0$.

\begin{thm}{Subsequences}{subseq}
Let $S_n\to S$ and let $n_1<n_2<n_3<\cdots$ be positive integers. Then $S_{n_j}\to S$ as $j\to\infty$. Consequently, if two subsequences of $\{S_n\}$ have different limits, then $\{S_n\}$ diverges.
\end{thm}
\begin{proof}
First, $n_j\ge j$ for every $j$: $n_1\ge1$, and if $n_j\ge j$ then $n_{j+1}>n_j\ge j$, so $n_{j+1}\ge j+1$. Let $\varepsilon>0$ and choose $N$ with $\abs{S_n-S}<\varepsilon$ for $n>N$. For $j>N$ we have $n_j\ge j>N$, so $\abs{S_{n_j}-S}<\varepsilon$. The second statement is the contrapositive combined with uniqueness of limits (Theorem~\ref{thm:unique}): if $S_n\to S$, every subsequence tends to the same $S$.
\end{proof}

We use the definition of the limit of a \emph{function} at infinity from Calculus~I: $\lim_{x\to\infty}g(x)=L$ means that for every $\varepsilon>0$ there is a number $X$ such that $\abs{g(x)-L}<\varepsilon$ for all $x>X$; and $\lim_{x\to\infty}g(x)=\infty$ means that for every $K$ there is $X$ with $g(x)>K$ for all $x>X$.

\begin{thm}{A monotone function has a limit at infinity}{monofn}
Let $g$ be defined on $[a,\infty)$ and nondecreasing, that is, $g(x)\le g(y)$ whenever $a\le x<y$.
\begin{enumerate}[label=(\alph*)]
  \item If $g$ is bounded above, then $\lim_{x\to\infty}g(x)$ exists and equals $\sup\{g(x):x\ge a\}$.
  \item If $g$ is not bounded above, then $g(x)\to\infty$ as $x\to\infty$.
\end{enumerate}
\end{thm}
\begin{proof}
(a) The set $\{g(x):x\ge a\}$ is nonempty and bounded above, so $M=\sup\{g(x):x\ge a\}$ exists by the Completeness Axiom. Let $\varepsilon>0$. By Theorem~\ref{thm:approx} there is $x_0\ge a$ with $M-\varepsilon<g(x_0)\le M$. For every $x>x_0$, monotonicity gives $M-\varepsilon<g(x_0)\le g(x)\le M$, so $\abs{g(x)-M}<\varepsilon$. Hence $g(x)\to M$.

(b) Let $K$ be a number. Since $K$ is not an upper bound of $g$, there is $x_0$ with $g(x_0)>K$, and then $g(x)\ge g(x_0)>K$ for all $x>x_0$.
\end{proof}

\begin{thm}{Improper integrals of nonnegative functions}{posint}
Let $f$ be continuous on $[a,\infty)$ with $f(x)\ge0$, and let $F(b)=\int_a^bf(x)\,dx$ for $b\ge a$. Then $F$ is nondecreasing, and:
\begin{enumerate}[label=(\alph*)]
  \item the improper integral $\int_a^\infty f(x)\,dx=\lim_{b\to\infty}F(b)$ converges if and only if $F$ is bounded above, and then $F(b)\le\int_a^\infty f(x)\,dx$ for every $b$;
  \item $F$ is bounded above if and only if the numbers $F(n)$, for integers $n\ge a$, are bounded above;
  \item if $F$ is not bounded above, then $\int_a^\infty f(x)\,dx=\infty$ (the integral diverges).
\end{enumerate}
\end{thm}
\begin{proof}
If $a\le b<c$ then $F(c)-F(b)=\int_b^cf(x)\,dx\ge0$ because $f\ge0$, so $F$ is nondecreasing. Parts (a) and (c) follow from Theorem~\ref{thm:monofn}. If $F$ is bounded above, the limit exists and equals $\sup F$, which is at least every $F(b)$. If $F$ is not bounded above, then $F(b)\to\infty$, so the improper integral does not converge; this proves (c) and shows that a convergent integral forces $F$ to be bounded above. For (b), one direction is trivial. Conversely, suppose $F(n)\le K$ for all integers $n\ge a$. Given any $b\ge a$, choose an integer $n\ge b$; then $F(b)\le F(n)\le K$.
\end{proof}

\begin{exmp}{Bounded or unbounded integrals}{posintex}
Decide whether $\int_1^\infty e^{-x}\,dx$ and $\int_1^\infty\frac1x\,dx$ converge, using Theorem~\ref{thm:posint}.

\textbf{Solution.} Both integrands are positive and continuous on $[1,\infty)$. For $F(b)=\int_1^be^{-x}\,dx=e^{-1}-e^{-b}$ we have $F(b)<e^{-1}$ for every $b$, so $F$ is bounded above and the integral converges, to $e^{-1}\approx0.368$.
%@ C-ex1a | simplify(integrate(exp(-x),(x,1,b)) - (exp(-1)-exp(-b))) == 0 and limit(exp(-1)-exp(-b), b, oo) == exp(-1) | F(b)=\int_1^be^{-x}\,dx=e^{-1}-e^{-b}
For $G(b)=\int_1^b\frac1x\,dx=\ln b$, the values $G(e^m)=m$ for $m=1,2,3,\dots$ are not bounded above, so by part (b) $G$ is not bounded above and the integral diverges. Notice that the integrand $\frac1x$ tends to $0$ and still the integral is infinite: ``small integrand'' does not mean ``finite integral'', exactly as ``small terms'' does not mean ``finite sum''.
%@ C-ex1b | integrate(1/x,(x,1,b)) == log(b) and limit(log(b), b, oo) == oo and [log(exp(i)) for i in range(1,5)] == [1,2,3,4] | $G(e^m)=m$ for $m=1,2,3,\dots$
\end{exmp}

\begin{exmp}{Subsequences settle a question about divergence}{subseqex}
Show that the sequence $S_n=\frac{1+(-1)^n}2$ (that is, $0,1,0,1,\dots$) diverges, and that $S_n=\frac{1-(-1)^n}2$ has the same property.

\textbf{Solution.} The even-numbered terms are $S_{2j}=\frac{1+1}2=1$ for every $j$, a constant sequence with limit $1$. The odd-numbered terms are $S_{2j-1}=\frac{1-1}2=0$, a constant sequence with limit $0$. These are two subsequences with different limits, so by Theorem~\ref{thm:subseq} the sequence diverges.
%@ C-ex2 | [(1+(-1)**(2*i))/Integer(2) for i in range(1,5)] == [1,1,1,1] and [(1+(-1)**(2*i-1))/Integer(2) for i in range(1,5)] == [0,0,0,0] | $S_{2j-1}=\frac{1-1}2=0$
The second sequence is the first one shifted by one place, with the roles of $0$ and $1$ swapped, and the same argument applies.
\end{exmp}

% =====================================================================
\section{Splitting Sequences, Sign Parts and Rearranging Series}

This section supports the sections on alternating series and on absolute and conditional convergence. The first two results are proved in full. Rearranging series is different: the two main theorems are \emph{stated and sketched, but their full proofs are omitted} here (they need more room than a reference appendix has). One rearrangement is worked out completely, because it shows why the theorems are needed.

\begin{thm}{Splitting a sequence into residue classes}{split}
Let $q\ge2$ be an integer and $\{S_n\}$ a sequence. If, for each $r=0,1,\dots,q-1$, the subsequence $S_{qj+r}$ (the terms whose index leaves remainder $r$ on division by $q$) tends to the same number $L$ as $j\to\infty$, then $S_n\to L$. In particular (take $q=2$): if $S_{2j}\to L$ and $S_{2j-1}\to L$, then $S_n\to L$.
\end{thm}
\begin{proof}
Let $\varepsilon>0$. For each $r$ there is a number $J_r$ with $\abs{S_{qj+r}-L}<\varepsilon$ for all $j>J_r$. Put $J=\max\{J_0,\dots,J_{q-1}\}$ and $N=q(J+1)$. Let $n>N$ and write $n=qj+r$ with $0\le r<q$. Then $qj+r>q(J+1)$ gives $j>J+1-\frac rq>J$ (because $r<q$), so $j>J\ge J_r$ and $\abs{S_n-L}<\varepsilon$. Hence $S_n\to L$. For $q=2$ the odd terms are $S_{2j-1}=S_{2(j-1)+1}$, a shift of the class $r=1$, and a shifted sequence has the same limit (Theorem~\ref{thm:tail}).
\end{proof}

\begin{exmp}{Even terms and odd terms}{splitex}
The sequence $S_n=\dfrac{n+(-1)^n}{n}$ has $S_{2j}=1+\dfrac1{2j}\to1$ and $S_{2j-1}=1-\dfrac1{2j-1}\to1$. Conclude that $S_n\to1$, and explain why one class alone would not be enough.

\textbf{Solution.} Both classes tend to $1$, so $S_n\to1$ by Theorem~\ref{thm:split} (here also $S_n=1+\frac{(-1)^n}n$ and the Squeeze Theorem gives the same result). One class alone is not enough: for $T_n=1+(-1)^n$, the even terms are all $2$, so $T_{2j}\to2$, but the odd terms are all $0$, so $T_n$ has no limit. Theorem~\ref{thm:split} needs \emph{every} class to have the \emph{same} limit.
%@ D-ex1 | all(simplify((2*j+(-1)**(2*j))/(2*j) - (1+Q(1,2*j))) == 0 for j in range(1,8)) and limit(1+1/(2*n),n,oo) == 1 and limit(1-1/(2*n-1),n,oo) == 1 and [1+(-1)**i for i in range(1,7)] == [0,2,0,2,0,2] | $T_{2j}\to2$
\end{exmp}

For a series $\sum a_k$ write $a_k^+=\max\{a_k,0\}$ (the positive part) and $a_k^-=\max\{-a_k,0\}$ (the negative part). Both are $\ge0$, and
\[
a_k=a_k^+-a_k^-,\qquad \abs{a_k}=a_k^++a_k^- ,\qquad a_k^+=\frac{\abs{a_k}+a_k}2,\qquad a_k^-=\frac{\abs{a_k}-a_k}2 .
\]

\begin{thm}{Positive and negative parts}{parts}
\begin{enumerate}[label=(\alph*)]
  \item If $\sum\abs{a_k}$ converges, then $\sum a_k^+$ and $\sum a_k^-$ both converge.
  \item If $\sum a_k$ converges but $\sum\abs{a_k}$ diverges, then $\sum a_k^+$ and $\sum a_k^-$ both diverge: their partial sums are not bounded above.
\end{enumerate}
\end{thm}
\begin{proof}
Let $P_n=\sum_{k=1}^na_k^+$, $M_n=\sum_{k=1}^na_k^-$, $A_n=\sum_{k=1}^n\abs{a_k}=P_n+M_n$ and $S_n=\sum_{k=1}^na_k=P_n-M_n$.

(a) The sequences $\{P_n\}$ and $\{M_n\}$ are nondecreasing (their terms are $\ge0$) and $P_n\le A_n\le A$, $M_n\le A_n\le A$, where $A=\sum\abs{a_k}$ is the limit of the nondecreasing sequence $\{A_n\}$ (Theorem~\ref{thm:mctbound}). By the Monotone Convergence Theorem both converge.

(b) Suppose $\{P_n\}$ converged, to $P$. Then $M_n=P_n-S_n\to P-S$, so $A_n=P_n+M_n$ would converge, contradicting the divergence of $\sum\abs{a_k}$. The same argument with the roles of $P_n$ and $M_n$ exchanged (use $P_n=M_n+S_n$) shows that $\{M_n\}$ does not converge either. A nondecreasing sequence that does not converge is not bounded above, by the Monotone Convergence Theorem.
\end{proof}

\begin{exmp}{The parts of the alternating harmonic series}{partsex}
For $a_k=\dfrac{(-1)^{k+1}}k$ find $a_k^+$ and $a_k^-$ and describe their partial sums.

\textbf{Solution.} $a_k^+=\frac1k$ for odd $k$ and $0$ for even $k$; $a_k^-=\frac1k$ for even $k$ and $0$ for odd $k$. So $\sum a_k^+=1+\frac13+\frac15+\cdots$ and $\sum a_k^-=\frac12+\frac14+\frac16+\cdots$. Since the alternating harmonic series converges but $\sum\frac1k$ diverges, Theorem~\ref{thm:parts}(b) says both of these diverge. This can also be seen directly: $\sum_{j=1}^m\frac1{2j-1}>\sum_{j=1}^m\frac1{2j}=\frac12H_m\to\infty$.
%@ D-ex2 | all(sum(Q(1,2*j-1) for j in range(1,m_+1)) > sum(Q(1,2*j) for j in range(1,m_+1)) for m_ in range(1,9)) and all(sum(Q(1,2*j) for j in range(1,m_+1)) == harmonic(m_)/2 for m_ in range(1,9)) | $\sum_{j=1}^m\frac1{2j-1}>\sum_{j=1}^m\frac1{2j}=\frac12H_m\to\infty$
\end{exmp}

\subsection*{Rearranging a series}

A \emph{rearrangement} of $\sum a_k$ is a series $\sum a_{\sigma(k)}$ that has the same terms in a different order ($\sigma$ is a one-to-one correspondence of the positive integers with themselves). For finite sums the order never matters. For infinite sums it can.

\begin{exmp}{Reordering the alternating harmonic series halves its sum}{rearr}
Let $S=1-\frac12+\frac13-\frac14+\cdots$ be the alternating harmonic series (it converges, by the Alternating Series Test treated in the section on alternating series; its value is $\ln2\approx0.693$, as the chapter on power series will show). Take its terms in this order: one positive term, then two negative terms,
\[
R=1-\frac12-\frac14+\frac13-\frac16-\frac18+\frac15-\frac1{10}-\frac1{12}+\cdots,
\]
so group $j$ is $\dfrac1{2j-1}-\dfrac1{4j-2}-\dfrac1{4j}$. Every term of $S$ appears exactly once in $R$: the positive terms $\frac1{2j-1}$ in order; the negative terms $-\frac1{4j-2}=-\frac1{2(2j-1)}$ and $-\frac1{4j}$ are the terms $-\frac1m$ with $m\equiv2$ and $m\equiv0\pmod4$, that is all the even $m$. Show that $R=\frac12S$. In particular $R\ne S$.

\textbf{Solution.} Let $R_n$ be the partial sums of $R$ and $S_n$ those of $S$. Since $\frac1{4j-2}=\frac12\cdot\frac1{2j-1}$, group $j$ equals
\[
\frac1{2j-1}-\frac1{2(2j-1)}-\frac1{4j}=\frac12\cdot\frac1{2j-1}-\frac12\cdot\frac1{2j}=\frac12\left(\frac1{2j-1}-\frac1{2j}\right).
\]
%@ D-rearr1 | all(simplify(Q(1,2*j-1)-Q(1,4*j-2)-Q(1,4*j) - Q(1,2)*(Q(1,2*j-1)-Q(1,2*j))) == 0 for j in range(1,10)) | \frac12\cdot\frac1{2j-1}-\frac12\cdot\frac1{2j}=\frac12\left(\frac1{2j-1}-\frac1{2j}\right).
Adding $j=1,\dots,m$ gives $R_{3m}=\frac12S_{2m}$. Because $S_{2m}\to S$, $R_{3m}\to\frac12S$. The other two classes differ from $R_{3m}$ by the next terms: $R_{3m+1}=R_{3m}+\frac1{2m+1}$ and $R_{3m+2}=R_{3m}+\frac1{2m+1}-\frac1{4m+2}$, and both added amounts tend to $0$. So all three subsequences $R_{3m}$, $R_{3m+1}$, $R_{3m+2}$ tend to $\frac12S$, and $R_n\to\frac12S$ by Theorem~\ref{thm:split}. Finally $S>0$: the partial sums $S_{2m}$ are nondecreasing (each step adds $\frac1{2m-1}-\frac1{2m}>0$), so $S\ge S_2=\frac12$ by Theorem~\ref{thm:mctbound}, and $R=\frac12S\ne S$.
%@ D-rearr2 | all(sum(Q(1,2*j-1)-Q(1,4*j-2)-Q(1,4*j) for j in range(1,m_+1)) == Q(1,2)*sum(Q(-1,i)*(-1)**i for i in range(1,2*m_+1)) for m_ in range(1,10)) | Adding $j=1,\dots,m$ gives $R_{3m}=\frac12S_{2m}$.
%@ D-rearr3 | all(sum(Q(1,2*j-1)-Q(1,4*j-2)-Q(1,4*j) for j in range(1,m_+1)) + Q(1,2*m_+1) - Q(1,4*m_+2) == sum(Q(1,2*j-1)-Q(1,4*j-2)-Q(1,4*j) for j in range(1,m_+2)) + Q(1,4*m_+4) for m_ in range(1,8)) | $R_{3m+2}=R_{3m}+\frac1{2m+1}-\frac1{4m+2}$
%@ D-rearr4 | abs(float(Q(1,2)*log(2))-0.3466)<1e-4 and abs(sum(float(Q(1,2*j-1)-Q(1,4*j-2)-Q(1,4*j)) for j in range(1,20001))-float(log(2)/2))<1e-4 and abs(sum(float(Q(1,i))*(-1)**(i+1) for i in range(1,20001))-float(log(2)))<1e-4 | $S\ge S_2=\frac12$
The same terms, in a different order, have a different sum.
\end{exmp}

Two theorems describe exactly what can happen. \textbf{The proofs below are omitted}; only a sketch is given.

\begin{thm}{Rearranging an absolutely convergent series}{rearrabs}
If $\sum\abs{a_k}$ converges, then every rearrangement of $\sum a_k$ converges, and to the same sum as $\sum a_k$.
\end{thm}
%%@omitted rearrabs
\emph{Proof omitted. Sketch.} Let $S$ be the sum and $\varepsilon>0$. Choose $N$ with $\sum_{k>N}\abs{a_k}<\varepsilon$. In the rearranged series, wait until the terms $a_1,\dots,a_N$ have all appeared, say by position $M$. For $n\ge M$ the partial sum of the rearrangement is $a_1+\cdots+a_N$ plus some terms $a_k$ with $k>N$, while the true partial sum $S_N$ is $a_1+\cdots+a_N$ alone. So the rearranged partial sum differs from $S_N$ by at most $\sum_{k>N}\abs{a_k}<\varepsilon$, and $S_N$ differs from $S$ by at most the same amount. The step that needs care is that $\abs{\sum_{k>N}a_k}\le\sum_{k>N}\abs{a_k}$ for the infinite tail.

\begin{thm}{Rearranging a conditionally convergent series (Riemann)}{riemann}
If $\sum a_k$ converges but $\sum\abs{a_k}$ diverges, then for every real number $A$ there is a rearrangement of $\sum a_k$ whose sum is $A$. There are also rearrangements whose partial sums tend to $\infty$, and rearrangements that have no limit.
\end{thm}
%%@omitted riemann
\emph{Proof omitted. Sketch.} By Theorem~\ref{thm:parts}(b) the positive terms add to $+\infty$ and the negative terms to $-\infty$, while $a_k\to0$. To aim at $A$: add positive terms, in their original order, until the running total first exceeds $A$; then add negative terms until it first drops below $A$; then positive terms again; and so on. This uses every term exactly once, because each stage needs only finitely many terms (the supply of each sign is unlimited). After the first crossing of $A$, the running total never misses $A$ by more than the size of the last term used, and the terms tend to $0$, so the running totals tend to $A$. The details of ``never misses by more than'' and of showing that every term is used are what the full proof supplies.

\begin{exmp}{Aiming at $\frac32$}{aim}
Follow the sketch for the alternating harmonic series with the target $A=\frac32$, and list the first stages.

\textbf{Solution.} The positive terms are $1,\frac13,\frac15,\dots$ and the negative terms are $-\frac12,-\frac14,\dots$ Add positive terms until the total exceeds $\frac32$: $1+\frac13=\frac43<\frac32$, and $\frac43+\frac15=\frac{23}{15}\approx1.533>\frac32$. Now add negative terms until the total drops below $\frac32$: already the first one does, $\frac{23}{15}-\frac12=\frac{31}{30}\approx1.033<\frac32$. Add positive terms again: $\frac{31}{30}+\frac17+\frac19+\frac1{11}+\frac1{13}\approx1.4551<\frac32$, and one more, $\frac1{15}$, gives $\approx1.5218>\frac32$. The series begins $1+\frac13+\frac15-\frac12+\frac17+\frac19+\frac1{11}+\frac1{13}+\frac1{15}-\frac14+\cdots$, and its totals swing around $\frac32$ by smaller and smaller amounts (a computer run of $10^5$ terms ends within $10^{-4}$ of $\frac32$).
%@ D-aim1 | Q(4,3) < Q(3,2) and Q(4,3)+Q(1,5) == Q(23,15) and Q(23,15) > Q(3,2) and Q(23,15)-Q(1,2) == Q(31,30) and Q(31,30) < Q(3,2) | $1+\frac13=\frac43<\frac32$, and $\frac43+\frac15=\frac{23}{15}\approx1.533>\frac32$
%@ D-aim2 | abs(float(Q(31,30)+Q(1,7)+Q(1,9)+Q(1,11)+Q(1,13))-1.4551)<1e-4 and Q(31,30)+Q(1,7)+Q(1,9)+Q(1,11)+Q(1,13) < Q(3,2) and abs(float(Q(31,30)+Q(1,7)+Q(1,9)+Q(1,11)+Q(1,13)+Q(1,15))-1.5218)<1e-4 | $\frac{31}{30}+\frac17+\frac19+\frac1{11}+\frac1{13}\approx1.4551<\frac32$
\end{exmp}

% =====================================================================
\section{Limits of $n$-th Roots}

The Root Test asks for the limit of $\abs{a_k}^{1/k}$, and for the common terms of series (powers of $k$ times a constant, or a rational function of $k$) that limit is $\lim n^{1/n}$ or one very close to it. This section proves the limits that are needed, using only the Squeeze Theorem (Theorem~\ref{thm:squeeze}), the limit laws (Theorem~\ref{thm:laws}), limits of $f(a_n)$ (Theorem~\ref{thm:fofa}) and the binomial theorem. One further limit is used in the section on the Ratio and Root Tests and is \emph{not} proved here: $\left(1+\frac1n\right)^n\to e$, which is how the number $e$ is defined in Calculus I.

\begin{thm}{The limit of $n^{1/n}$}{nthroot}
$\displaystyle\lim_{n\to\infty}n^{1/n}=1$.
\end{thm}
\begin{proof}
Let $n\ge2$ and write $n^{1/n}=1+h_n$. Since $n\ge1$, we have $n^{1/n}\ge1$, so $h_n\ge0$. By the binomial theorem, all of whose terms are nonnegative here,
\[
n=(1+h_n)^n=1+nh_n+\binom n2h_n^2+\cdots\ \ge\ \binom n2h_n^2=\frac{n(n-1)}2\,h_n^2 .
\]
Dividing by $\frac{n(n-1)}2>0$ gives $h_n^2\le\frac2{n-1}$, so $0\le h_n\le\sqrt{2/(n-1)}$, that is,
\[
1\le n^{1/n}\le1+\sqrt{\frac2{n-1}} .
\]
%@ E-nr | all(float(n_)**(1/float(n_)) >= 1 and float(n_)**(1/float(n_)) <= 1+(2/(n_-1))**0.5 and n_ >= Q(n_*(n_-1),2)*(float(n_)**(1/float(n_))-1)**2 - 1e-9 for n_ in range(2,300)) | 1\le n^{1/n}\le1+\sqrt{\frac2{n-1}} .
The right side tends to $1$ (limit laws, and Theorem~\ref{thm:fofa}(a) for the square root, which is continuous at $0$). By the Squeeze Theorem, $n^{1/n}\to1$.
\end{proof}

\begin{thm}{Roots of constants}{croot}
If $c>0$, then $\displaystyle\lim_{n\to\infty}c^{1/n}=1$.
\end{thm}
\begin{proof}
Suppose first that $c\ge1$. For every $n\ge c$ we have $1\le c\le n$, and $t\mapsto t^{1/n}$ is increasing, so $1\le c^{1/n}\le n^{1/n}$. Both outer sequences tend to $1$ (Theorem~\ref{thm:nthroot}), so $c^{1/n}\to1$ by the Squeeze Theorem. Now suppose $0<c<1$. Then $\frac1c>1$, so $\left(\frac1c\right)^{1/n}\to1$ by the first case, and
\[
c^{1/n}=\frac1{(1/c)^{1/n}}\to\frac11=1
\]
by the limit laws.
\end{proof}

\begin{thm}{Roots of powers times a positive limit}{posroot}
Let $p$ be a real number and let $c_n>0$ with $c_n\to L$, where $L>0$. Then $\displaystyle\lim_{n\to\infty}\left(n^pc_n\right)^{1/n}=1$.
\end{thm}
\begin{proof}
\emph{The factor $c_n$.} Take $\varepsilon=\frac L2$ in the definition of the limit: there is $N$ with $\frac L2<c_n<\frac{3L}2$ for $n>N$. Raising to the power $\frac1n$ (an increasing operation), $\left(\frac L2\right)^{1/n}<c_n^{1/n}<\left(\frac{3L}2\right)^{1/n}$. Both outer sequences tend to $1$ (Theorem~\ref{thm:croot}), so $c_n^{1/n}\to1$ by the Squeeze Theorem.

\emph{The factor $n^p$.} $(n^p)^{1/n}=\left(n^{1/n}\right)^p\to1^p=1$, by Theorem~\ref{thm:nthroot} and Theorem~\ref{thm:fofa}(a) applied to $y\mapsto y^p$, which is continuous at $1$.

\emph{Product.} For positive numbers $(xy)^{1/n}=x^{1/n}y^{1/n}$, so $\left(n^pc_n\right)^{1/n}=\left(n^{1/n}\right)^p\,c_n^{1/n}\to1\cdot1=1$ by the product law.
\end{proof}

\begin{exmp}{How slowly $n^{1/n}$ approaches $1$}{nrtable}
Compute $n^{1/n}$ for $n=2,3,10,100,1000$ and compare with the bound in the proof.

\textbf{Solution.} $2^{1/2}\approx1.4142$, $3^{1/3}\approx1.4423$, $10^{1/10}\approx1.2589$, $100^{1/100}\approx1.0471$, $1000^{1/1000}\approx1.0069$. The sequence first \emph{rises} (to its largest value at $n=3$, because $\frac{\ln x}x$ has its maximum at $x=e$) and only then falls to $1$, and it falls slowly. The bound $1+\sqrt{2/(n-1)}$ is valid but crude: at $n=1000$ it is $\approx1.0447$, while the true value is $1.0069$. So the limit is a statement about very large $n$.
%@ E-ex1 | abs(float(2)**0.5-1.4142)<1e-4 and abs(3**(1/3)-1.4423)<1e-4 and abs(10**0.1-1.2589)<1e-4 and abs(100**0.01-1.0471)<1e-4 and abs(1000**0.001-1.0069)<1e-4 and all(3**(1/3) >= j**(1/j) for j in range(1,500)) and abs(1+(2/999)**0.5-1.0447)<1e-4 | $10^{1/10}\approx1.2589$
\end{exmp}

\begin{exmp}{Roots of rational expressions}{nrex}
Find (a)~$\displaystyle\lim_{n\to\infty}\left(\frac{3n^2+1}{n^2+5}\right)^{1/n}$, (b)~$\displaystyle\lim_{n\to\infty}\left(7n^{10}\right)^{1/n}$, (c)~$\displaystyle\lim_{n\to\infty}\left(\frac{5n+1}{n^4}\right)^{1/n}$.

\textbf{Solution.} (a) The fraction tends to $3>0$, so with $p=0$ Theorem~\ref{thm:posroot} gives $1$. (b) Take $p=10$ and $c_n=7$: the limit is $1$. (c) Write $\frac{5n+1}{n^4}=n^{-3}\left(5+\frac1n\right)$: here $p=-3$ and $c_n=5+\frac1n\to5>0$, so the limit is $1$. In every case the root of a quantity that is only polynomial in $n$ is $1$ in the limit, and only a genuinely exponential factor ($2^n$, $\left(1+\frac1n\right)^{n^2}$, \dots) changes the answer.
%@ E-ex2 | abs(float(((3*1000**2+1)/(1000**2+5))**(1/1000))-1.0011)<1e-4 and abs(float((7*1000**10)**(1/1000))-1.0736)<1e-3 and limit(5+1/n,n,oo) == 5 and simplify((5*n+1)/n**4 - n**(-3)*(5+1/n)) == 0 and limit((3*n**2+1)/(n**2+5),n,oo) == 3 | \frac{5n+1}{n^4}=n^{-3}\left(5+\frac1n\right)
\end{exmp}

% =====================================================================
\section{Power Series: the Radius of Convergence and Term-by-Term Calculus}

Chapter 11 uses three facts about a power series $\sum c_k(x-a)^k$ that need more care than a lecture allows. First, it converges on an interval around its center, so that a radius of convergence exists. Second, inside that interval it can be differentiated and integrated term by term. Third, doing so does not change the radius. This section proves all three \emph{in full}. It uses the Divergence Test, the Comparison Test, the Ratio Test, the geometric series and the Mean Value Theorem, and (for the radius) the least upper bound property of Section~A. Throughout, $a$ is the center and $c_k$ are the coefficients.

\begin{thm}{Convergence at one point controls all nearer points}{abel}
If $\sum c_k(x_0-a)^k$ converges, then $\sum c_k(x-a)^k$ converges absolutely for every $x$ with $\abs{x-a}<\abs{x_0-a}$.
\end{thm}
\begin{proof}
Put $s=\abs{x_0-a}$. If $s=0$ no $x$ satisfies $\abs{x-a}<s$ and there is nothing to prove, so let $s>0$. The terms of a convergent series tend to $0$ (otherwise it would diverge, by the Divergence Test), and a convergent sequence is bounded (Theorem~\ref{thm:bounded}). So there is a number $B$ with $\abs{c_k}s^k\le B$ for every $k$. Let $\abs{x-a}<s$ and put $t=\abs{x-a}/s$, so $0\le t<1$. Then
\[
\abs{c_k(x-a)^k}=\abs{c_k}s^k\,t^k\le Bt^k .
\]
The geometric series $\sum Bt^k$ converges because $t<1$, so $\sum\abs{c_k(x-a)^k}$ converges by the Comparison Test.
\end{proof}

\begin{thm}{The radius of convergence exists}{radius}
A power series $\sum c_k(x-a)^k$ converges in exactly one of three ways: (1)~it converges absolutely for every $x$ (then $R=\infty$); (2)~there is a number $R>0$ such that it converges absolutely for $\abs{x-a}<R$ and diverges for $\abs{x-a}>R$; (3)~it converges only at $x=a$ (then $R=0$). The number $R$ is determined by the series.
\end{thm}
\begin{proof}
Let $S$ be the set of numbers $\abs{x-a}$ for the $x$ at which the series converges. It contains $0$ (take $x=a$), so $S$ is not empty. There are three possibilities for $S$.

\emph{$S$ is not bounded above.} Let $x$ be any real number. Some element of $S$ exceeds $\abs{x-a}$, say $\abs{x_0-a}$ with the series convergent at $x_0$. By Theorem~\ref{thm:abel} the series converges absolutely at $x$. This is case~(1).

\emph{$S=\{0\}$.} The series converges only at $a$: case~(3).

\emph{$S$ is bounded above and has a positive element.} By the Completeness Axiom (Section~A), $R=\sup S$ exists, and $R>0$. Let $\abs{x-a}<R$. Then $\abs{x-a}$ is not an upper bound of $S$ (because $R$ is the \emph{least} upper bound), so some $s\in S$ has $s>\abs{x-a}$; the series converges at a point at distance $s$ from $a$, and by Theorem~\ref{thm:abel} it converges absolutely at $x$. Let $\abs{x-a}>R$. Every element of $S$ is $\le R$, so $\abs{x-a}\notin S$: the series diverges at $x$. This is case~(2).

The three cases are exhaustive and exclusive, and $R$ is determined: if two numbers $R<R'$ both satisfied~(2), then at a point with $R<\abs{x-a}<R'$ the series would both diverge and converge.
\end{proof}

\begin{thm}{The differentiated and integrated series have the same radius}{radderiv}
Let $\sum c_k(x-a)^k$ have radius $R$. For $\abs{x-a}<R$ the three series
\[
\sum_{k=1}^\infty kc_k(x-a)^{k-1},\qquad\sum_{k=2}^\infty k(k-1)c_k(x-a)^{k-2},\qquad\sum_{k=0}^\infty\frac{c_k}{k+1}(x-a)^{k+1}
\]
all converge absolutely. The first and the third have radius exactly $R$.
\end{thm}
\begin{proof}
\emph{Convergence.} Fix $x$ and put $\rho=\abs{x-a}<R$. Choose $s$ with $\rho<s<R$. By Theorem~\ref{thm:radius} the original series converges at the point $a+s$, so its terms tend to $0$ and $\abs{c_k}s^k\le B$ for some $B$ (Theorem~\ref{thm:bounded}). With $t=\rho/s\in[0,1)$,
\[
k\abs{c_k}\rho^{k-1}\le\frac Bs\,k\,t^{k-1},\qquad k(k-1)\abs{c_k}\rho^{k-2}\le\frac B{s^2}\,k(k-1)\,t^{k-2}.
\]
The series $\sum kt^{k-1}$ and $\sum k(k-1)t^{k-2}$ converge for $0\le t<1$: for $t=0$ they have one nonzero term at most, and for $0<t<1$ the Ratio Test applies, since the ratios $\frac{(k+1)t}k$ and $\frac{(k+1)t}{k-1}$ tend to $t<1$. By the Comparison Test the first two series converge absolutely. For the third, $\frac{\abs{c_k}\rho^{k+1}}{k+1}\le\rho\abs{c_k}\rho^k$, and $\sum\abs{c_k}\rho^k$ converges because $\rho<R$ (Theorem~\ref{thm:radius}); compare again.

\emph{The radii.} Let $R_D$ and $R_G$ be the radii of the first and third series. We have shown $R_D\ge R$ and $R_G\ge R$. Conversely, let $\abs u<R_D$ (with $u$ measured from the center $a$). The first series converges absolutely at $u$, and $\abs{c_k}\abs u^k\le k\abs{c_k}\abs u^k=\abs u\cdot k\abs{c_k}\abs u^{k-1}$ for $k\ge1$, so $\sum\abs{c_k}\abs u^k$ converges. Then the original series converges at $u$, so $\abs u\le R$ (it diverges beyond $R$). Hence $R_D\le R$. For the third series, observe that its \emph{differentiated} series is the original one, because $k\cdot\frac{c_{k-1}}k=c_{k-1}$. Applying what we showed ($R_D\ge R$) to the third series, the original converges absolutely wherever $\abs u<R_G$, so $R\ge R_G$. Therefore $R_D=R_G=R$.
\end{proof}

\begin{thm}{Differentiating a power series term by term}{termdiff}
Let $\sum c_k(x-a)^k$ have radius $R>0$ and $f(x)=\sum_{k=0}^\infty c_k(x-a)^k$ for $\abs{x-a}<R$. Then $f$ is differentiable at every $x$ with $\abs{x-a}<R$, and $f'(x)=\sum_{k=1}^\infty kc_k(x-a)^{k-1}$.
\end{thm}
\begin{proof}
Fix $x_0$ with $\abs{x_0-a}<R$ and choose $\rho$ with $\abs{x_0-a}<\rho<R$. Put $u=x_0-a$ and, for $h\ne0$ so small that $\abs{u+h}<\rho$, $v=u+h$. All three series $f(x_0+h)=\sum c_kv^k$, $f(x_0)=\sum c_ku^k$ and $\sum kc_ku^{k-1}$ converge (Theorem~\ref{thm:radderiv}), and so does $E=\sum_{k\ge2}k(k-1)\abs{c_k}\rho^{k-2}$, by the same theorem at the point at distance $\rho$. By the limit laws for series,
\[
\frac{f(x_0+h)-f(x_0)}h-\sum_{k=1}^\infty kc_ku^{k-1}=\sum_{k=2}^\infty c_k\left(\frac{v^k-u^k}h-ku^{k-1}\right),
\]
because the terms $k=0$ and $k=1$ of the right side are $0$. Fix $k\ge2$. By the Mean Value Theorem for $t\mapsto t^k$ on the interval between $u$ and $v$, $v^k-u^k=k\xi^{k-1}h$ for some $\xi$ between $u$ and $v$. Applying it again to $t\mapsto t^{k-1}$ on the interval between $u$ and $\xi$ gives $\xi^{k-1}-u^{k-1}=(k-1)\eta^{k-2}(\xi-u)$ with $\eta$ between $u$ and $\xi$. Hence
\[
\frac{v^k-u^k}h-ku^{k-1}=k\left(\xi^{k-1}-u^{k-1}\right)=k(k-1)\eta^{k-2}(\xi-u),
\]
and since $\eta$ and $\xi$ lie between $u$ and $v$ (so $\abs\eta<\rho$) and $\abs{\xi-u}\le\abs h$, the absolute value of the $k$-th term of the right side is at most $\abs{c_k}\,k(k-1)\rho^{k-2}\abs h$. The partial sums of the right side are therefore bounded in absolute value by $E\abs h$, and by Theorem~\ref{thm:ineq} so is the sum. So
\[
\abs{\frac{f(x_0+h)-f(x_0)}h-\sum_{k=1}^\infty kc_ku^{k-1}}\le E\,\abs h\longrightarrow0\quad(h\to0),
\]
that is, $f'(x_0)=\sum kc_ku^{k-1}$.
\end{proof}

\begin{thm}{Integrating a power series term by term}{termint}
With $R$ and $f$ as in Theorem~\ref{thm:termdiff}, the series $G(x)=\sum_{k=0}^\infty\frac{c_k}{k+1}(x-a)^{k+1}$ converges for $\abs{x-a}<R$ and $G'(x)=f(x)$ there. Consequently $\int f(x)\,dx=G(x)+C$ on $(a-R,a+R)$.
\end{thm}
\begin{proof}
By Theorem~\ref{thm:radderiv}, $G$ is a power series centered at $a$ with radius $R$. By Theorem~\ref{thm:termdiff} applied to $G$, $G'(x)=\sum_{k=0}^\infty(k+1)\frac{c_k}{k+1}(x-a)^k=f(x)$ for $\abs{x-a}<R$. Two antiderivatives of $f$ on an interval differ by a constant (a consequence of the Mean Value Theorem), so every antiderivative is $G+C$.
\end{proof}

\begin{thm}{Repeated differentiation}{allderiv}
With $R$ and $f$ as in Theorem~\ref{thm:termdiff}, $f$ has derivatives of every order on $(a-R,a+R)$, and for $j\ge0$
\[
f^{(j)}(x)=\sum_{k=j}^\infty k(k-1)\cdots(k-j+1)\,c_k(x-a)^{k-j},\qquad f^{(j)}(a)=j!\,c_j .
\]
\end{thm}
\begin{proof}
For $j=0$ this is the definition of $f$. If the formula holds for $j$, its right side is a power series centered at $a$ with the coefficients $k(k-1)\cdots(k-j+1)c_k$, and by Theorem~\ref{thm:radderiv} (applied $j$ times) it has radius $R$. Theorem~\ref{thm:termdiff} applied to it gives the formula for $j+1$. At $x=a$ every term with $k>j$ contains $(x-a)^{k-j}=0$, and the term $k=j$ is $j!\,c_j$.
\end{proof}

\begin{exmp}{The key estimate for $\frac1{1-x}$}{termdiffex}
Take $f(x)=\frac1{1-x}=\sum x^k$ ($a=0$, $R=1$), $x_0=0.3$, $h=0.01$ and $\rho=\frac12$. Compare the actual error in the difference quotient with the bound $E\abs h$ in the proof of Theorem~\ref{thm:termdiff}.

\textbf{Solution.} $f'(x_0)=\frac1{0.7^2}=\frac{100}{49}\approx2.040816$. The difference quotient is $\frac{f(0.31)-f(0.3)}{0.01}=100\left(\frac{100}{69}-\frac{10}7\right)=\frac{1000}{483}\approx2.070393$, so the actual difference is $\frac{700}{23667}\approx0.0296$. The bound uses $E=\sum_{k\ge2}k(k-1)\rho^{k-2}=\frac2{(1-\rho)^3}=16$ (the second derivative of $\frac1{1-x}$ at $\rho$), so $E\abs h=0.16$. The bound is valid and five times larger than the true difference, as is typical: it must work for every $c_k$ and every $h$.
%@ F-ex1 | Q(100,69)-Q(10,7) == Q(10,483) and 100*Q(10,483) == Q(1000,483) and Q(1000,483)-Q(100,49) == Q(700,23667) and abs(float(Q(700,23667))-0.0296)<1e-4 and summation(k*(k-1)*Q(1,2)**(k-2),(k,2,oo)) == 16 and 16*Q(1,100) == Q(4,25) | $E=\sum_{k\ge2}k(k-1)\rho^{k-2}=\frac2{(1-\rho)^3}=16$
\end{exmp}

% =====================================================================
\section{Abel's Theorem and the Binomial Series at the Endpoints}

The section on Taylor series needs two facts about the \emph{endpoints} of an interval of convergence, where the theorems on power series say nothing. First, a power series can converge at an endpoint $x=a+R$ and equal a function $f$ inside the interval; is its sum at the endpoint the number $f(a+R)$? Abel's theorem says that the sum is the \emph{limit} of $f$ as $x$ approaches the endpoint from inside, which is $f(a+R)$ when $f$ is continuous there. This is how $1-\frac12+\frac13-\cdots=\ln2$ and $1-\frac13+\frac15-\cdots=\frac\pi4$ are proved. Second, the binomial series $\sum\binom pkx^k$ has radius $1$, and whether it converges at $x=1$ and at $x=-1$ depends on $p$; we settle every case. This section proves both facts \emph{in full}. It uses the Divergence Test, the Alternating Series Test, the Comparison Test, the harmonic series, the geometric series and the Monotone Convergence Theorem (Section~A).

\begin{thm}{Abel's limit theorem}{abellimit}
Let $\sum c_k(x-a)^k$ have radius of convergence $R$ with $0<R<\infty$, and let $f(x)=\sum_{k=0}^\infty c_k(x-a)^k$ for $\abs{x-a}<R$.
\begin{enumerate}
  \item If $\sum_{k=0}^\infty c_kR^k$ converges, then $\displaystyle\lim_{x\to(a+R)^-}f(x)=\sum_{k=0}^\infty c_kR^k$.
  \item If $\sum_{k=0}^\infty c_k(-R)^k$ converges, then $\displaystyle\lim_{x\to(a-R)^+}f(x)=\sum_{k=0}^\infty c_k(-R)^k$.
\end{enumerate}
\end{thm}
\begin{proof}
(1)~\emph{Reduction.} Put $d_k=c_kR^k$ and $t=\frac{x-a}R$. Then $f(x)=g(t)$ with $g(t)=\sum d_kt^k$ for $\abs t<1$, and $x\to(a+R)^-$ means $t\to1^-$. Let $S=\sum d_k$. Replacing $d_0$ by $d_0-S$ subtracts the constant $S$ from $g$ and makes $\sum d_k=0$. So it is enough to prove: \emph{if $\sum d_k$ converges to $0$, then $g(t)\to0$ as $t\to1^-$.}

\emph{Summation by parts.} Let $s_n=d_0+\cdots+d_n$, so $s_n\to0$, and put $s_{-1}=0$, so $d_k=s_k-s_{k-1}$. For every $n\ge1$ and every $t$,
\[
\sum_{k=0}^nd_kt^k=\sum_{k=0}^ns_kt^k-\sum_{k=0}^ns_{k-1}t^k=\sum_{k=0}^ns_kt^k-\sum_{k=0}^{n-1}s_kt^{k+1}=(1-t)\sum_{k=0}^{n-1}s_kt^k+s_nt^n .
\]
Now let $\abs t<1$. The sequence $(s_n)$ converges, so it is bounded, $\abs{s_k}\le K$, and $\sum s_kt^k$ converges absolutely by comparison with the geometric series $\sum Kt^k$; also $s_nt^n\to0$. Letting $n\to\infty$,
\[
g(t)=(1-t)\sum_{k=0}^\infty s_kt^k\qquad(\abs t<1).
\]

\emph{The estimate.} Let $\varepsilon>0$. Choose $N$ with $\abs{s_k}<\varepsilon$ for all $k>N$, and put $M=\sum_{k=0}^N\abs{s_k}$. For $0<t<1$,
\[
\abs{g(t)}\le(1-t)\sum_{k=0}^N\abs{s_k}t^k+(1-t)\sum_{k=N+1}^\infty\varepsilon\,t^k\le(1-t)M+\varepsilon\,t^{N+1}\le(1-t)M+\varepsilon ,
\]
because $(1-t)\sum_{k>N}t^k=t^{N+1}\le1$. If $1-t<\frac\varepsilon{M+1}$, then $\abs{g(t)}<2\varepsilon$. Since $\varepsilon$ was arbitrary, $g(t)\to0$ as $t\to1^-$.

(2)~Put $\tilde c_k=(-1)^kc_k$. For $\abs{x-a}<R$ we have $\sum\tilde c_k(x-a)^k=\sum c_k(a-x)^k=f(2a-x)$, and $\sum\tilde c_kR^k=\sum c_k(-R)^k$ converges. By part~(1), $f(2a-x)\to\sum c_k(-R)^k$ as $x\to(a+R)^-$. Writing $y=2a-x$, which tends to $(a-R)^+$, gives part~(2).
\end{proof}

\begin{exmp}{The hypothesis cannot be dropped}{abelfail}
The series $\sum_{k=0}^\infty(-1)^kx^k$ has radius $1$ and sums to $f(x)=\frac1{1+x}$ for $\abs x<1$. As $x\to1^-$, $f(x)\to\frac12$. Is $\frac12$ the ``sum of the series at $x=1$''?

\textbf{Solution.} No: at $x=1$ the series is $1-1+1-1+\cdots$, whose terms $(-1)^k$ do not tend to $0$, so it diverges (Divergence Test), and there is no sum for the limit to equal. Numerically, $f(0.9)=\frac1{1.9}\approx0.526$ and $f(0.99)=\frac1{1.99}\approx0.5025$ approach $\frac12$. Abel's theorem applies only when the series converges at the endpoint; it then says that the limit of $f$ \emph{is} the sum.
%@ G-abelfail | limit(1/(1+x),x,1,'-') == Q(1,2) and abs(float(1/(1+Q(9,10)))-0.526)<1e-3 and abs(float(1/(1+Q(99,100)))-0.5025)<1e-4 | \frac1{1.99}\approx0.5025
\end{exmp}

\begin{exmp}{Abel's theorem and the alternating harmonic series}{abelln2}
Show that $\sum_{k=1}^\infty\frac{(-1)^{k+1}}k=\ln2$, given that $\ln(1+x)=\sum_{k=1}^\infty\frac{(-1)^{k+1}x^k}k$ for $\abs x<1$.

\textbf{Solution.} The power series has radius $1$ and converges at $x=1$ by the Alternating Series Test (the terms $\frac1k$ decrease to $0$). By Theorem~\ref{thm:abellimit}(1), its sum at $x=1$ equals $\lim_{x\to1^-}\ln(1+x)=\ln2$, because the logarithm is continuous at $2$. The convergence is slow: the error after $n$ terms is at most $\frac1{n+1}$, and the first ten terms give $0.6456$ against $\ln2=0.6931$.
%@ G-abelln2 | limit(log(1+x),x,1,'-') == log(2) and summation((-1)**(k+1)/k,(k,1,oo)) == log(2) and abs(float(sum(Integer(-1)**(j+1)*Q(1,j) for j in range(1,11)))-0.6456)<1e-4 and abs(float(log(2))-0.6931)<1e-4 | $0.6456$ against $\ln2=0.6931$
\end{exmp}

For the second fact, recall that for a real number $p$ and an integer $k\ge0$ the \emph{binomial coefficient} is $\binom pk=\frac{p(p-1)\cdots(p-k+1)}{k!}$, with $\binom p0=1$. If $p$ is a nonnegative integer then $\binom pk=0$ for $k>p$ (the factor $p-p=0$ appears) and $\sum\binom pkx^k$ is a polynomial. Otherwise no factor is $0$ and every $\binom pk\ne0$. Let $c_k=\binom pk$ and $u_k=\abs{c_k}$. Multiplying out the definition,
\[
(k+1)\,c_{k+1}=(p-k)\,c_k,\qquad\text{so}\qquad(k+1)\,u_{k+1}=\abs{p-k}\,u_k .\tag{$\ast$}
\]
We also use $1-y\le e^{-y}$ for all real $y$ (the graph of the convex function $e^{-y}$ lies above its tangent line $1-y$ at $y=0$).

\begin{thm}{The size of the binomial coefficients}{binomsize}
Let $p$ be a real number that is not a nonnegative integer.
\begin{enumerate}
  \item If $p\le-1$ then $u_k\ge1$ for every $k$.
  \item If $-1<p<0$ then $c_k=(-1)^ku_k$, the numbers $u_k$ decrease strictly to $0$, and $\sum u_k$ diverges.
  \item If $p>0$, let $K$ be the smallest integer greater than $p$. For $k\ge K$ the signs of $c_k$ alternate, the numbers $u_k$ decrease strictly to $0$, and $\sum u_k$ converges.
\end{enumerate}
\end{thm}
\begin{proof}
(1)~Here $\abs{p-k}=k+\abs p\ge k+1$, so $(\ast)$ gives $u_{k+1}\ge u_k$. Since $u_0=1$, $u_k\ge1$ for all $k$.

(2)~Every factor $p-j$ is negative, so $c_k$ has the sign $(-1)^k$. Now $\abs{p-k}=k+\abs p$ with $0<\abs p<1$, so $(\ast)$ reads $(k+1)u_{k+1}=(k+\abs p)u_k$, which gives three facts. First, $u_{k+1}<u_k$ because $k+\abs p<k+1$. Second, $(k+1)u_{k+1}>k\,u_k$, so $k\,u_k$ increases and $k\,u_k\ge1\cdot u_1=\abs p$ for $k\ge1$; that is $u_k\ge\frac{\abs p}k$, and $\sum u_k$ diverges by comparison with the harmonic series. Third, $\frac{u_{k+1}}{u_k}=1-\frac{1-\abs p}{k+1}\le e^{-(1-\abs p)/(k+1)}$, so
\[
0<u_n=u_0\prod_{j=0}^{n-1}\frac{u_{j+1}}{u_j}\le\exp\left(-(1-\abs p)\sum_{j=1}^n\frac1j\right)\to0,
\]
because the harmonic series diverges. So $u_n\to0$ by the Squeeze Theorem.

(3)~For $k\ge K$ we have $k>p$, so $p-k<0$. Then $\frac{c_{k+1}}{c_k}=\frac{p-k}{k+1}<0$: the signs alternate. By $(\ast)$, $(k+1)u_{k+1}=(k-p)u_k$, so $u_{k+1}<u_k$. As in~(2), $\frac{u_{k+1}}{u_k}=1-\frac{p+1}{k+1}\le e^{-(p+1)/(k+1)}$, so for $n>K$, $0<u_n\le u_K\exp\left(-(p+1)\sum_{j=K+1}^n\frac1j\right)\to0$. Finally $(\ast)$ gives $k\,u_k-(k+1)u_{k+1}=p\,u_k$ for $k\ge K$. Adding these equations for $k=K,\dots,n$, the left sides telescope:
\[
p\sum_{k=K}^nu_k=K\,u_K-(n+1)\,u_{n+1}\le K\,u_K .
\]
So the partial sums of $\sum u_k$ are bounded, and since the terms are positive the partial sums increase, so $\sum u_k$ converges by the Monotone Convergence Theorem (Theorem~\ref{thm:mct}).
\end{proof}

\begin{thm}{The binomial series at the endpoints}{binomend}
Let $p$ be a real number that is not a nonnegative integer. The series $\sum_{k=0}^\infty\binom pkx^k$ has radius of convergence $1$, and
\begin{enumerate}
  \item at $x=1$ it converges if and only if $p>-1$;
  \item at $x=-1$ it converges if and only if $p>0$.
\end{enumerate}
So its interval of convergence is $(-1,1)$ if $p\le-1$, $(-1,1]$ if $-1<p<0$, and $[-1,1]$ if $p>0$.
\end{thm}
\begin{proof}
\emph{Radius.} $\abs{c_{k+1}/c_k}=\frac{\abs{p-k}}{k+1}\to1$, so $R=1$ (the Ratio Test applied to the coefficients).

\emph{At $x=1$.} If $p\le-1$ the terms satisfy $\abs{c_k}\ge1$ (Theorem~\ref{thm:binomsize}(1)), so they do not tend to $0$ and the series diverges (Divergence Test). If $-1<p<0$ then $c_k=(-1)^ku_k$ with $u_k$ decreasing to $0$ (Theorem~\ref{thm:binomsize}(2)), so the series converges by the Alternating Series Test. If $p>0$, the terms from $k=K$ on alternate in sign and decrease in absolute value to $0$ (Theorem~\ref{thm:binomsize}(3)); the Alternating Series Test applied to the tail shows that the series converges (the first $K$ terms are finitely many).

\emph{At $x=-1$.} The terms are $(-1)^kc_k$. If $p\le-1$ their absolute values are $\ge1$: divergence. If $-1<p<0$ they are $(-1)^kc_k=u_k>0$, and $\sum u_k$ diverges (Theorem~\ref{thm:binomsize}(2)). If $p>0$, $\sum u_k$ converges (Theorem~\ref{thm:binomsize}(3); the first $K$ terms do not matter), so the series converges absolutely.
\end{proof}

\begin{exmp}{The coefficients for $p=\frac12$}{binomsum}
For $p=\frac12$ we have $K=1$ in Theorem~\ref{thm:binomsize}(3), and the telescoping identity reads $\frac12\sum_{k=1}^nu_k=u_1-(n+1)u_{n+1}$, that is
\[
\sum_{k=1}^n\abs{\binom{1/2}k}=1-2(n+1)\abs{\binom{1/2}{n+1}} .
\]
Check it for $n=9$ and say what it shows about $\sum_{k=1}^\infty\abs{\binom{1/2}k}$.

\textbf{Solution.} The absolute values of the coefficients are $u_1=\frac12$, $u_2=\frac18$, $u_3=\frac1{16}$, $u_4=\frac5{128}$, $u_5=\frac7{256}$, $u_6=\frac{21}{1024}$, $u_7=\frac{33}{2048}$, $u_8=\frac{429}{32768}$, $u_9=\frac{715}{65536}$, $u_{10}=\frac{2431}{262144}$. Their sum for $k=1,\dots,9$ is $\frac{53381}{65536}\approx0.81453$, and the right side is $1-20\cdot\frac{2431}{262144}=\frac{53381}{65536}$: equal. The partial sums increase and are at most $1$, so they converge to some $L\le1$, and then $(n+1)u_{n+1}=\frac{1-S_n}2\to\frac{1-L}2$. If this limit were positive, $u_n\ge\frac cn$ for large $n$ and $\sum u_k$ would diverge by comparison with the harmonic series, contradicting Theorem~\ref{thm:binomsize}(3). So the limit is $0$, $L=1$, and $\sum_{k=1}^\infty\abs{\binom{1/2}k}=1$. Since $(-1)^k\binom{1/2}k=-u_k$ for $k\ge1$, this says that $\sum_{k=0}^\infty\binom{1/2}k(-1)^k=1-1=0=\sqrt{1+(-1)}$.
%@ G-binomsum | [abs(binomial(Q(1,2),j)) for j in range(1,11)] == [Q(1,2),Q(1,8),Q(1,16),Q(5,128),Q(7,256),Q(21,1024),Q(33,2048),Q(429,32768),Q(715,65536),Q(2431,262144)] and sum(abs(binomial(Q(1,2),j)) for j in range(1,10)) == Q(53381,65536) and 1-20*Q(2431,262144) == Q(53381,65536) and abs(float(Q(53381,65536))-0.81453)<1e-5 and all(sum(abs(binomial(Q(1,2),j)) for j in range(1,nn+1)) == 1-2*(nn+1)*abs(binomial(Q(1,2),nn+1)) for nn in range(1,25)) | $\frac{53381}{65536}\approx0.81453$
\end{exmp}

% =====================================================================
\section{Multiplying and Dividing Power Series}

Section 11.4 multiplies and divides Taylor series to find the first terms of functions such as $e^x\sin x$ and $\tan x$, and uses the product of two series to prove that $e^xe^y=e^{x+y}$. The product of two sums is \emph{not} the sum of the products of corresponding terms, so this needs a proof: the rule is the Cauchy product, in which the coefficient of $(x-a)^n$ collects every product $a_jb_k$ with $j+k=n$. This section proves that the rule is correct (Theorems~\ref{thm:cauchy} and~\ref{thm:psprod}), that a power series with nonzero constant term has a reciprocal that is again a power series (Theorem~\ref{thm:recip}), and hence that quotients work too (Theorem~\ref{thm:quot}). Everything is proved in full. It uses the Monotone Convergence Theorem (Section~A), the existence of the radius of convergence (Section~F) and the uniqueness of power series coefficients (Theorem~\ref{thm:allderiv}).

\begin{thm}{Products of absolutely convergent series}{cauchy}
Let $\sum_{j=0}^\infty u_j$ and $\sum_{j=0}^\infty v_j$ be absolutely convergent series with sums $U$ and $V$, and let $w_n=\sum_{j=0}^nu_jv_{n-j}$. Then $\sum_{n=0}^\infty w_n$ converges absolutely and its sum is $UV$.
\end{thm}
\begin{proof}
Put $\mu=\sum\abs{u_j}$, $\nu=\sum\abs{v_j}$, $U_N=\sum_{i=0}^Nu_i$, $V_N=\sum_{j=0}^Nv_j$ and $S_N=\sum_{n=0}^Nw_n$. Then $S_N$ is the sum of $u_iv_j$ over the pairs of integers $i,j\ge0$ with $i+j\le N$, and $U_NV_N$ is the sum over the pairs with $i,j\le N$.

\emph{Absolute convergence.} The pairs with $i+j\le N$ all have $i,j\le N$, so
\[
\sum_{n=0}^N\abs{w_n}\le\sum_{i+j\le N}\abs{u_i}\abs{v_j}\le\Bigl(\sum_{i=0}^N\abs{u_i}\Bigr)\Bigl(\sum_{j=0}^N\abs{v_j}\Bigr)\le\mu\nu .
\]
The partial sums of $\sum\abs{w_n}$ are nondecreasing and bounded above by $\mu\nu$, so the series converges (Theorem~\ref{thm:mct}).

\emph{The sum.} The difference $U_NV_N-S_N$ is the sum of $u_iv_j$ over the pairs with $i,j\le N$ and $i+j>N$. If both $i\le N/2$ and $j\le N/2$ then $i+j\le N$, so in each such pair $i>N/2$ or $j>N/2$. Let $M=\lfloor N/2\rfloor$ and let $\mu_M=\sum_{i>M}\abs{u_i}$ and $\nu_M=\sum_{j>M}\abs{v_j}$ be the tails; they tend to $0$ as $N\to\infty$ (a tail is $\mu$ minus a partial sum that tends to $\mu$). Splitting the pairs according to which of $i>M$ and $j>M$ holds (a pair for which both hold is counted twice, which only enlarges the bound),
\[
\abs{U_NV_N-S_N}\le\sum_{i>M}\abs{u_i}\sum_{j\ge0}\abs{v_j}+\sum_{i\ge0}\abs{u_i}\sum_{j>M}\abs{v_j}\le\mu_M\,\nu+\mu\,\nu_M\longrightarrow0 .
\]
Since $U_NV_N\to UV$ (Theorem~\ref{thm:laws}), $S_N=U_NV_N-(U_NV_N-S_N)\to UV$.
\end{proof}

\begin{thm}{Products of power series}{psprod}
Let $f(x)=\sum_{k=0}^\infty a_k(x-a)^k$ and $g(x)=\sum_{k=0}^\infty b_k(x-a)^k$ for all $x$ with $\abs{x-a}<R$, where $R>0$, and put $c_n=\sum_{j=0}^na_jb_{n-j}$. Then $\sum_{n=0}^\infty c_n(x-a)^n$ converges to $f(x)g(x)$ for every $x$ with $\abs{x-a}<R$.
\end{thm}
\begin{proof}
Let $\abs{x-a}<R$ and choose $x_1$ with $\abs{x-a}<\abs{x_1-a}<R$. Both series converge at $x_1$, so by Theorem~\ref{thm:abel} both converge absolutely at $x$. Apply Theorem~\ref{thm:cauchy} to $u_j=a_j(x-a)^j$ and $v_j=b_j(x-a)^j$: then $w_n=\sum_{j=0}^na_jb_{n-j}(x-a)^{j}(x-a)^{n-j}=c_n(x-a)^n$, and the theorem gives $\sum c_n(x-a)^n=f(x)g(x)$.
\end{proof}

\begin{thm}{The reciprocal of a power series}{recip}
Let $f(x)=\sum_{k=0}^\infty a_k(x-a)^k$ for $\abs{x-a}<R$, where $R>0$ and $a_0\ne0$. Define $b_0=\frac1{a_0}$ and $b_n=-\frac1{a_0}\sum_{j=1}^na_jb_{n-j}$ for $n\ge1$. Then there is $\rho>0$ such that for $\abs{x-a}<\rho$ the series $\sum b_n(x-a)^n$ converges, $f(x)\ne0$ and $\dfrac1{f(x)}=\sum_{n=0}^\infty b_n(x-a)^n$.
\end{thm}
\begin{proof}
\emph{The coefficients do not grow fast.} Choose $s$ with $0<s<R$. The series $\sum a_ks^k$ converges, so its terms tend to $0$ and are bounded: $\abs{a_k}s^k\le B$ for every $k$, where we may take $B\ge\abs{a_0}$. Put $L=B/\abs{a_0}$ and $d_i=s^i\abs{b_i}$, so $d_0=\frac1{\abs{a_0}}$. For $n\ge1$ the definition of $b_n$ gives
\[
\abs{b_n}\le\frac1{\abs{a_0}}\sum_{j=1}^n\abs{a_j}\abs{b_{n-j}}\le L\sum_{j=1}^ns^{-j}\abs{b_{n-j}},\qquad\text{so}\qquad d_n\le L\sum_{i=0}^{n-1}d_i .
\]
We claim $d_n\le d_0(1+L)^n$ for all $n$. It is true for $n=0$; if it holds for all indices below $n\ge1$, then $d_n\le L\,d_0\sum_{i=0}^{n-1}(1+L)^i=d_0\bigl((1+L)^n-1\bigr)\le d_0(1+L)^n$. Hence $\abs{b_n}\le d_0\left(\frac{1+L}s\right)^n$.

\emph{Convergence.} Let $\rho=\frac s{1+L}$ (so $\rho<s<R$) and $\abs{x-a}<\rho$. Then $\abs{b_n(x-a)^n}\le d_0\theta^n$ with $\theta=\frac{(1+L)\abs{x-a}}s<1$, so $\sum b_n(x-a)^n$ converges (comparison with a geometric series); call its sum $g(x)$.

\emph{The product is $1$.} By Theorem~\ref{thm:psprod} (with $\rho$ in place of $R$), $f(x)g(x)=\sum c_n(x-a)^n$ with $c_n=\sum_{j=0}^na_jb_{n-j}$. Here $c_0=a_0b_0=1$, and for $n\ge1$, $c_n=a_0b_n+\sum_{j=1}^na_jb_{n-j}=0$ by the definition of $b_n$. So $f(x)g(x)=1$ for $\abs{x-a}<\rho$; therefore $f(x)\ne0$ and $g(x)=\frac1{f(x)}$.
\end{proof}

\begin{thm}{Quotients of power series}{quot}
Let $f(x)=\sum a_k(x-a)^k$ and $g(x)=\sum b_k(x-a)^k$ for $\abs{x-a}<R$, with $R>0$ and $b_0\ne0$. Then there is $\rho>0$ such that $\dfrac{f(x)}{g(x)}=\sum_{n=0}^\infty q_n(x-a)^n$ for $\abs{x-a}<\rho$, where the coefficients are determined one after another by
\[
q_n=\frac1{b_0}\Bigl(a_n-\sum_{j=1}^nb_jq_{n-j}\Bigr),\qquad\text{equivalently}\qquad a_n=\sum_{j=0}^nb_jq_{n-j}.
\]
\end{thm}
\begin{proof}
By Theorem~\ref{thm:recip}, $\frac1g=\sum r_n(x-a)^n$ for $\abs{x-a}<\rho_1$, for some $\rho_1>0$ and coefficients $r_n$. Then $\frac fg=f\cdot\frac1g=\sum q_n(x-a)^n$ on $\abs{x-a}<\min(\rho_1,R)=:\rho$ by Theorem~\ref{thm:psprod}, where $q_n=\sum_{j=0}^na_jr_{n-j}$. Apply Theorem~\ref{thm:psprod} again to $g$ and $\sum q_n(x-a)^n$: for $\abs{x-a}<\rho$,
\[
f(x)=g(x)\cdot\frac{f(x)}{g(x)}=\sum_{n=0}^\infty\Bigl(\sum_{j=0}^nb_jq_{n-j}\Bigr)(x-a)^n .
\]
Both sides are power series in $(x-a)$ equal on an interval, so their coefficients agree (by Theorem~\ref{thm:allderiv}, the coefficient of $(x-a)^n$ in a power series is $\frac{1}{n!}$ times the $n$-th derivative of its sum at $a$). Hence $a_n=\sum_{j=0}^nb_jq_{n-j}$; solving for $q_n$ (the term $j=0$ is $b_0q_n$) gives the formula.
\end{proof}

\begin{exmp}{The product $e^x\cdot\frac1{1-x}$}{prodex}
Find the first terms of the Maclaurin series of $\dfrac{e^x}{1-x}$, and the limit of its coefficients.

\textbf{Solution.} With $a_j=\frac1{j!}$ (valid for all $x$) and $b_k=1$ (valid for $\abs x<1$), Theorem~\ref{thm:psprod} with $R=1$ gives $c_n=\sum_{j=0}^n\frac1{j!}\cdot1$, the partial sums of the series for $e$:
\[
\frac{e^x}{1-x}=1+2x+\frac52x^2+\frac83x^3+\frac{65}{24}x^4+\frac{163}{60}x^5+\cdots\qquad(\abs x<1).
\]
%@ H-prod | Tay(exp(x)/(1-x),0,5) == 1+2*x+Q(5,2)*x**2+Q(8,3)*x**3+Q(65,24)*x**4+Q(163,60)*x**5 and [sum(Q(1,factorial(j)) for j in range(nn+1)) for nn in range(6)] == [1,2,Q(5,2),Q(8,3),Q(65,24),Q(163,60)] | +\frac{65}{24}x^4+\frac{163}{60}x^5+\cdots
The coefficients increase to $e=2.71828\dots$, because $c_n=\sum_{j\le n}\frac1{j!}\to e$.
\end{exmp}

\begin{exmp}{The first terms of $\tan x$}{tanex}
Use Theorem~\ref{thm:quot} with $f=\sin x$ and $g=\cos x$ to find the first four nonzero terms of the series of $\tan x$.

\textbf{Solution.} $f$ has $a_1=1$, $a_3=-\frac16$, $a_5=\frac1{120}$, $a_7=-\frac1{5040}$ and zero for the other indices up to $7$; $g$ has $b_0=1$, $b_2=-\frac12$, $b_4=\frac1{24}$, $b_6=-\frac1{720}$ and zero for odd indices. By the formula $q_n=a_n-\sum_{j=1}^nb_jq_{n-j}$ (here $b_0=1$): $q_0=0$; $q_1=a_1=1$; $q_2=-b_2q_0=0$; $q_3=a_3-b_2q_1=-\frac16+\frac12=\frac13$; $q_4=0$; $q_5=a_5-b_2q_3-b_4q_1=\frac1{120}+\frac16-\frac1{24}=\frac2{15}$; $q_7=a_7-b_2q_5-b_4q_3-b_6q_1=-\frac1{5040}+\frac1{15}-\frac1{72}+\frac1{720}=\frac{17}{315}$. So
\[
\tan x=x+\frac{x^3}3+\frac{2x^5}{15}+\frac{17x^7}{315}+\cdots
\]
on some interval about $0$ (Theorem~\ref{thm:quot} guarantees one; its length is not found by this method).
%@ H-tan | Tay(tan(x),0,7) == x+x**3/3+2*x**5/15+17*x**7/315 and -Q(1,6)+Q(1,2) == Q(1,3) and Q(1,120)+Q(1,6)-Q(1,24) == Q(2,15) and -Q(1,5040)+Q(1,15)-Q(1,72)+Q(1,720) == Q(17,315) | x+\frac{x^3}3+\frac{2x^5}{15}+\frac{17x^7}{315}+\cdots
\end{exmp}


\end{document}
